\chapter{Einführung}
\label{cha:Intro}

Beim Anfertigen einer Abschlussarbeit steht man als Studierender meist zum ersten Mal vor dem Problem, einen längeren wissenschaftlichen Text mit Bildern, Gleichungen und Referenzen schreiben zu müssen.
Dafür bietet sich das Textsatz-System \LaTeX\ an, zu dessen Vorteilen die weitgehende Trennung von Inhalt und Layout gehören.

Die Arbeiten sollen im Corporate Design der TU Darmstadt abgefasst werden.
Hierzu stellt die TU Darmstadt das \LaTeX-Paket \verb|tuda-ci| zur Verfügung, welches auch für den vorliegenden Text verwendet wurde.


Allgemeines zu den Grundlagen von \LaTeX{} findet man \zB in \cite{Kopka} oder \cite{Schmidt}.
Allgemeines über wissenschaftliche Arbeiten findet man in \cite{Friedrich}.
Als Einstieg in die Typografie ist \cite{Willberg} sehr zu empfehlen.

Die Arbeit kann in Deutsch oder wahlweise auch in Englisch verfasst werden.

In der vorliegenden Anleitung werden in \charef{cha:Hinweise} zuerst allgemeine Hinweise und Tipps zur Durchführung einer wissenschaftlichen Arbeit gegeben.
Im Anschluß daran geht \charef{cha:Hinweise_Latex} kurz auf die Verwendung von \LaTeX{} zum Erstellen der schriftlichen Arbeit ein.
In \charef{cha:Hinweise_Allgemein} wird dann wieder allgemein auf die Abfassung schriftlicher Arbeiten eingegangen.
Dazu wird zu jedem Punkt nach den allgemeinen Hinweisen auch die Umsetzung in \LaTeX{} angesprochen.

\charef{cha:Beurteilung} umreißt die Kriterien, nach denen eine Arbeit an unserem Fachgebiet bewertet wird.

Im Anhang findet sich eine Checkliste die vor Abgabe der Arbeit unbedingt abgearbeitet werden sollte, um zu prüfen, ob alle Vorgaben und Hinweise beachtet wurden.

Dieses Dokument ist \emph{keine} Einführung in \LaTeX{}.
Es werden zwar einige Befehle und Umgebungen namentlich genannt, jedoch sind diese Aufzählungen weder vollständig noch werden die Befehle hier genauer erläutert.
Hierzu sollten sich im Internet schnell alle Informationen finden lassen.
Erfahrungsgemäß ist es auch nicht nötig, spezielle Tutorials oder Kurse zu \LaTeX\ durchzuführen.
Ausgehend von dem zur Verfügung gestellten Minimalbeispiel lässt sich das notwendige Wissen in der Regel gut iterativ, während der Bearbeitung der Arbeit, erlangen.
Empfehlenswert ist es dazu, möglichst früh mit \LaTeX{} zu beginnen, und beispielsweise schon Stichworte in der Einarbeitungsphase der Arbeit zu "`texen"'.
